% BHU PYQ -- Transcribed & Corrected
% Paper: ENGAE-31: Reading Skills
% Exam: B.A. / B.Sc. / B.Com. (Hons.) Semester III (End Term) Examination, 2025-26 (NEP)
% Category: Ability Enhancement Course (AEC)

\documentclass[11pt,a4paper]{article}
\usepackage[a4paper,top=2.5cm,bottom=2.5cm,left=2.8cm,right=2.8cm,headheight=14pt]{geometry}
\usepackage{amsmath,amssymb}
\usepackage[T1]{fontenc}
\usepackage[utf8]{inputenc}
\usepackage{lmodern,microtype,fancyhdr,enumitem,hyperref,lastpage,parskip}

\hypersetup{
    colorlinks=true,
    urlcolor=black,
    linkcolor=black,
    pdftitle={ENGAE-31: Reading Skills -- BHU Sem III 2025-26 (NEP)}
}

\pagestyle{fancy}
\fancyhf{}
\renewcommand{\headrulewidth}{0.4pt}
\renewcommand{\footrulewidth}{0.4pt}
\fancyhead[L]{\small   \textit{Banaras Hindu University}}
\fancyhead[R]{\small   \textit{ENGAE-31: Reading Skills (AEC)}}
\fancyfoot[C]{\small Page  \thepage\ of \pageref{LastPage}}


\newlist{parts}{enumerate}{2}
\setlist[parts,1]{label=(\Alph*),leftmargin=*,itemsep=3pt,topsep=2pt}

\newcommand{\pts}[1]{\hfill   \textit{[#1]}}


\begin{document}


\begin{center}
  {\large\bfseries BANARAS HINDU UNIVERSITY}\\[2pt]
  {
\normalsize B.A. / B.Sc. / B.Com. (Hons.) --- Semester III Examination, 2025-26 (NEP)}\\[6pt]
  \rule{0.8\linewidth}{0.5pt}\\[6pt]
  {\large\bfseries Ability Enhancement Course (AEC)}\\[4pt]
  {\large\bfseries Paper Code: ENGAE-31 : Reading Skills}\\[4pt]
  {
\normalsize Time: 1.5 Hours\hfill Maximum Marks: 70}
\end{center}

\rule{\linewidth}{0.4pt}

\smallskip



\noindent   \textit{Instructions: Answer all questions. Each question carries equal marks (1 mark each).}

\bigskip




\noindent   \textbf{Question 1.} The best correction of the sentence, ``Their going to the market,'' would be:\pts{1}

\begin{parts}
  \item There going to the market
  \item Their gone to the market
  \item They're going to the market
  \item Its going to the market
\end{parts}

\medskip




\noindent   \textbf{Question 2.} The line ``The road was a ribbon of moonlight'' (from the poem   \textit{The Highwayman} by Alfred Noyes) is an example of:\pts{1}

\begin{parts}
  \item Literal meaning
  \item Metaphoric comparison
  \item Factual reporting
  \item Note-taking and summary
\end{parts}

\medskip




\noindent   \textbf{Question 3.} Writing the words ``cat,'' ``dog,'' and ``run'' as /kæt/, /dɒɡ/, and /rʌn/ is an example of:\pts{1}

\begin{parts}
  \item Phonic pronunciation
  \item Intensive reading
  \item Skimming a text
  \item Silent reading
\end{parts}

\medskip




\noindent   \textbf{Question 4.} Which of these is NOT an external barrier to a student's reading comprehension?\pts{1}

\begin{parts}
  \item Noisy classroom environment
  \item Lack of encouragement by teacher
  \item Jokes about pronunciation by other students
  \item Severe headache and fever
\end{parts}

\medskip




\noindent   \textbf{Question 5.} Sunil doesn't know the meaning of the word ``dubious,'' but is able to guess it from the sentence, ``I could never trust him and found all his excuses to be dubious.'' What is this an example of?\pts{1}

\begin{parts}
  \item Predictive reading / Contextual guessing
  \item Error detection and editing
  \item Extensive reading
  \item Graphic organization
\end{parts}

\medskip




\noindent   \textbf{Question 6.} Skimming a text is primarily used for:\pts{1}

\begin{parts}
  \item Finding the main idea or gist quickly
  \item Memorizing line by line
  \item Checking punctuation errors
  \item Translating word for word
\end{parts}

\medskip




\noindent   \textbf{Question 7.} Scanning a reading passage involves:\pts{1}

\begin{parts}
  \item Searching for a specific piece of information or fact
  \item Reading every sentence in detail
  \item Summarizing the whole chapter
  \item Pronouncing words aloud
\end{parts}

\medskip




\noindent   \textbf{Question 8.} Intensive reading focuses on:\pts{1}

\begin{parts}
  \item Detailed comprehension and careful analysis of text
  \item Fast casual reading for fun
  \item Skipping difficult words
  \item Looking at pictures only
\end{parts}

\medskip




\noindent   \textbf{Question 9.} Extensive reading encourages students to:\pts{1}

\begin{parts}
  \item Read large volumes of text for general understanding and pleasure
  \item Analyze every grammatical rule
  \item Memorize the entire book
  \item Avoid reading long stories
\end{parts}

\medskip




\noindent   \textbf{Question 10.} Inferential reading requires a reader to:\pts{1}

\begin{parts}
  \item Read between the lines to uncover hidden meanings
  \item Copy sentences verbatim
  \item Read only the titles
  \item Count the number of paragraphs
\end{parts}

\medskip




\noindent   \textbf{Question 11.} Critical reading involves:\pts{1}

\begin{parts}
  \item Evaluating the author's arguments, logic, and evidence
  \item Accepting all statements without questioning
  \item Reading as fast as possible
  \item Ignoring the author's perspective
\end{parts}

\medskip




\noindent   \textbf{Question 12.} Active reading strategies include:\pts{1}

\begin{parts}
  \item Annotating, highlighting, and asking questions while reading
  \item Passively skimming without taking notes
  \item Staring at text without focus
  \item Reading while watching television
\end{parts}

\medskip




\noindent   \textbf{Question 13.} A reader uses background knowledge (schema) to:\pts{1}

\begin{parts}
  \item Connect prior experience with new text information
  \item Forget previous lessons
  \item Ignore the author's topic
  \item Replace reading with guessing
\end{parts}

\medskip




\noindent   \textbf{Question 14.} Identifying the central thesis of an essay helps to understand:\pts{1}

\begin{parts}
  \item The main point or main argument of the author
  \item The font size used by the publisher
  \item The cost of the book
  \item The biographical background of the reader
\end{parts}

\medskip




\noindent   \textbf{Question 15.} In reading, context clues help a student to:\pts{1}

\begin{parts}
  \item Determine the meaning of unfamiliar words from surrounding text
  \item Skip the page entirely
  \item Memorize dictionary definitions
  \item Avoid reading long paragraphs
\end{parts}

\medskip




\noindent   \textbf{Question 16.} A summary of a reading passage should be:\pts{1}

\begin{parts}
  \item Brief, objective, and focused on key points
  \item Longer than the original text
  \item Full of personal opinions and slang
  \item An exact word-for-word copy
\end{parts}

\medskip




\noindent   \textbf{Question 17.} Paraphrasing means:\pts{1}

\begin{parts}
  \item Restating an author's ideas in your own words
  \item Quoting directly with quotation marks
  \item Changing only one word in a sentence
  \item Deleting paragraphs
\end{parts}

\medskip




\noindent   \textbf{Question 18.} Figurative language in literature includes:\pts{1}

\begin{parts}
  \item Similes, metaphors, and personification
  \item Literal dictionary definitions only
  \item Numerical statistics
  \item Bibliographical footnotes
\end{parts}

\medskip




\noindent   \textbf{Question 19.} Reading comprehension is best defined as:\pts{1}

\begin{parts}
  \item Understanding, interpreting, and making sense of written text
  \item Pronouncing words correctly without knowing their meaning
  \item Memorizing sentence structures
  \item Speed reading without retention
\end{parts}

\medskip




\noindent   \textbf{Question 20.} SQ3R study method stands for:\pts{1}

\begin{parts}
  \item Survey, Question, Read, Recite, Review
  \item Speak, Quote, Read, Write, Repeat
  \item Scan, Quick, Read, Rewrite, Record
  \item Search, Question, Reasoning, Review, Report
\end{parts}

\medskip




\noindent   \textbf{Question 21.} An author's purpose in writing a text can be to:\pts{1}

\begin{parts}
  \item Inform, persuade, or entertain
  \item Confuse the audience deliberately
  \item Waste paper
  \item Print random symbols
\end{parts}

\medskip




\noindent   \textbf{Question 22.} The tone of a text reflects:\pts{1}

\begin{parts}
  \item The author's attitude toward the subject or audience
  \item The size of the paper
  \item The length of the chapter
  \item The time of day it was read
\end{parts}

\medskip




\noindent   \textbf{Question 23.} Which reading speed is best suited for reviewing textbook material before an exam?\pts{1}

\begin{parts}
  \item Flexible reading speed adjusted to difficulty
  \item Ultra-fast skimming without stopping
  \item Slow word-by-word decoding
  \item Reading backwards
\end{parts}

\medskip




\noindent   \textbf{Question 24.} Graphic organizers like Mind Maps help readers to:\pts{1}

\begin{parts}
  \item Visualize relationships between ideas
  \item Avoid taking written notes
  \item Replace reading comprehension
  \item Color pages for decoration
\end{parts}

\medskip




\noindent   \textbf{Question 25.} Fact vs. Opinion: Which statement is a FACT?\pts{1}

\begin{parts}
  \item The sun rises in the east
  \item English is the most beautiful language
  \item Reading is more fun than sports
  \item Summer is better than winter
\end{parts}

\medskip




\noindent   \textbf{Question 26.} Fact vs. Opinion: Which statement is an OPINION?\pts{1}

\begin{parts}
  \item Water freezes at 0 degrees Celsius
  \item This novel is the most exciting book ever written
  \item Banaras Hindu University was founded in 1916
  \item Oxygen is required for human respiration
\end{parts}

\medskip




\noindent   \textbf{Question 27.} Making predictions while reading involves:\pts{1}

\begin{parts}
  \item Anticipating what will happen next based on text clues
  \item Guessing the price of the book
  \item Writing the ending before starting
  \item Reading only the back cover
\end{parts}

\medskip




\noindent   \textbf{Question 28.} Vocabulary in context can be understood using:\pts{1}

\begin{parts}
  \item Synonyms, antonyms, and explanation clues in the sentence
  \item Random guessing
  \item Skipping all difficult words
  \item Counting letters in a word
\end{parts}

\medskip




\noindent   \textbf{Question 29.} Which technique improves reading fluency?\pts{1}

\begin{parts}
  \item Reading in meaningful word chunks rather than word-by-word
  \item Finger pointing to every single letter
  \item Stopping after every word
  \item Reading aloud at maximum speed
\end{parts}

\medskip




\noindent   \textbf{Question 30.} Subvocalization refers to:\pts{1}

\begin{parts}
  \item Silent internal speech while reading text
  \item Shouting words in class
  \item Listening to an audiobook
  \item Writing notes on paper
\end{parts}

\medskip




\noindent   \textbf{Question 31.} A main idea is different from a supporting detail because:\pts{1}

\begin{parts}
  \item The main idea gives the overarching point, while details provide evidence
  \item Details are always placed in title headers
  \item Main ideas only appear in footnotes
  \item They are identical in function
\end{parts}

\medskip




\noindent   \textbf{Question 32.} When evaluating an online source for research reading, check for:\pts{1}

\begin{parts}
  \item Authority, accuracy, objectivity, and currency
  \item Bright colors and background music
  \item Number of pop-up advertisements
  \item Social media shares only
\end{parts}

\medskip




\noindent   \textbf{Question 33.} Common spelling errors due to homophones include confusing:\pts{1}

\begin{parts}
  \item ``Their'', ``There'', and ``They're''
  \item ``Cat'' and ``Dog''
  \item ``Book'' and ``Pen''
  \item ``Run'' and ``Walk''
\end{parts}

\medskip




\noindent   \textbf{Question 34.} Phonetic confusion in reading often happens when words:\pts{1}

\begin{parts}
  \item Sound similar but have different spellings and meanings
  \item Are written in uppercase letters
  \item Have short definitions
  \item Are printed in bold font
\end{parts}

\medskip




\noindent   \textbf{Question 35.} Confusing words like ``receive'' and ``relieve'' involves:\pts{1}

\begin{parts}
  \item Spelling rule exceptions (e.g.   \textit{i before e except after c})
  \item Math calculation errors
  \item Punctuation mistakes
  \item Translation errors
\end{parts}

\medskip




\noindent   \textbf{Question 36.} Engaged readers are typically:\pts{1}

\begin{parts}
  \item Interested, motivated, and persistent despite text challenges
  \item Passive and easily bored
  \item Unwilling to seek clarification
  \item Focused only on exam grades
\end{parts}

\medskip




\noindent   \textbf{Question 37.} Which of the following is NOT a true characteristic of reading?\pts{1}

\begin{parts}
  \item It does not improve overall communication skills
  \item It is an active cognitive process
  \item It can be both individual and collaborative
  \item It involves mental processing and comprehension
\end{parts}

\medskip




\noindent   \textbf{Question 38.} In group discussions about reading, students should focus on:\pts{1}

\begin{parts}
  \item Sharing ideas and experiences constructively
  \item Debating to win arguments
  \item Trying to read faster than others
  \item Criticizing classmates' opinions
\end{parts}

\medskip




\noindent   \textbf{Question 39.} A supportive classroom environment for reading development includes:\pts{1}

\begin{parts}
  \item Social, personal, and cognitive dimensions
  \item Strict silence and no discussion
  \item Rote memorization drills
  \item Speed tests only
\end{parts}

\medskip




\noindent   \textbf{Question 40.} Which of the following is NOT a core language domain?\pts{1}

\begin{parts}
  \item Scrolling social media reels
  \item Writing
  \item Listening and Speaking
  \item Reading
\end{parts}

\medskip




\noindent   \textbf{Question 41.} What is Reading Stamina?\pts{1}

\begin{parts}
  \item The ability to read long texts for extended periods of time with focus
  \item The speed at which you flip pages
  \item The volume at which you read aloud
  \item The ability to make dramatic expressions while reading
\end{parts}

\medskip




\noindent   \textbf{Question 42.} What does the step ``Breaking it Down'' signify in the reading process?\pts{1}

\begin{parts}
  \item Breaking a long, complex text into smaller, manageable pieces
  \item Reading aloud at high volume
  \item Breaking a promise to a friend
  \item Tearing pages out of a book
\end{parts}

\medskip




\noindent   \textbf{Question 43.} How does graphic organization help readers?\pts{1}

\begin{parts}
  \item It helps keep track of concepts and relationships visually
  \item It decorates notes with random pictures
  \item It prevents students from reading books
  \item It tests spelling speed
\end{parts}

\medskip




\noindent   \textbf{Question 44.} Setting reading goals is important for:\pts{1}

\begin{parts}
  \item Varying reading strategies according to specific purpose
  \item Memorizing entire texts word for word
  \item Avoiding texts altogether
  \item Reading in an unsystematic manner
\end{parts}

\medskip




\noindent   \textbf{Question 45.} Financial ``bank'' and river ``bank'' are examples of homonyms. Homonyms are defined as:\pts{1}

\begin{parts}
  \item Two words with the same spelling and pronunciation but different meanings
  \item Words with opposite meanings
  \item Words that rhyme at the end
  \item Words with different spellings and opposite meanings
\end{parts}

\medskip




\noindent   \textbf{Question 46.} What is the correct definition of Homophones?\pts{1}

\begin{parts}
  \item Two words with the same pronunciation but different spellings and meanings
  \item Two words with the exact same meaning and spelling
  \item Two words with opposite meanings
  \item Two words with different spellings and identical meanings
\end{parts}

\medskip




\noindent   \textbf{Question 47.} ``Hot'' and ``cold'' are antonyms. Antonyms are defined as:\pts{1}

\begin{parts}
  \item Two different words with directly opposing meanings
  \item Words with identical meanings
  \item Words that sound the same
  \item Words spelled identically
\end{parts}

\medskip




\noindent   \textbf{Question 48.} Which of the following is a pair of rhyming words?\pts{1}

\begin{parts}
  \item light, night
  \item off, rough
  \item tough, though
  \item move, love
\end{parts}

\medskip




\noindent   \textbf{Question 49.} Which of the following is NOT a pair of rhyming words?\pts{1}

\begin{parts}
  \item shove, groove
  \item cat, hat
  \item sing, ring
  \item book, look
\end{parts}

\medskip




\noindent   \textbf{Question 50.} In the sentence ``She can sings very well,'' the grammatical error is in:\pts{1}

\begin{parts}
  \item Verb form (``sings'')
  \item Subject (``She'')
  \item Word order
  \item Article
\end{parts}

\medskip




\noindent   \textbf{Question 51.} The sentence ``They goes to school every day'' contains an error in:\pts{1}

\begin{parts}
  \item Subject-verb agreement (``goes'')
  \item Word choice (``school'')
  \item Punctuation
  \item Preposition (``to'')
\end{parts}

\medskip




\noindent   \textbf{Question 52.} Which activity is used by teachers to prepare students before reading new text?\pts{1}

\begin{parts}
  \item Brainstorming: discussing prior knowledge about the topic
  \item Height-order seating arrangement
  \item Roll call sorting
  \item Silent exam grading
\end{parts}

\medskip




\noindent   \textbf{Question 53.} What is Visualization during reading?\pts{1}

\begin{parts}
  \item Being able to mentally picture scenes and information described in text
  \item Reading printed paper instead of a screen
  \item Staring at words without thinking
  \item Reading aloud off a blackboard
\end{parts}

\medskip




\noindent   \textbf{Question 54.} Which of the following is a positive peer interaction practice in reading class?\pts{1}

\begin{parts}
  \item Learning from each other's diverse backgrounds and perspectives
  \item Working on personal tasks while peers present
  \item Isolating students by region
  \item Avoiding group discussions
\end{parts}

\medskip




\noindent   \textbf{Question 55.} Fluency in English is valuable in India and globally because:\pts{1}

\begin{parts}
  \item It acts as a connecting bridge language for education and careers
  \item English is inherently superior to other languages
  \item Exams are impossible without it
  \item It prevents peer criticism
\end{parts}

\medskip




\noindent   \textbf{Question 56.} Which of the following is NOT a valid set of terms related to the reading process?\pts{1}

\begin{parts}
  \item singing, competition, drama
  \item text, engagement, competence
  \item cognition, reading strategy, purpose
  \item fluency, chunking, paraphrasing
\end{parts}

\medskip




\noindent   \textbf{Question 57.} Which of the following factors affects Readability of a text?\pts{1}

\begin{parts}
  \item All of the above (Vocabulary difficulty, sentence length, and concept density)
  \item Sentence length and vocabulary difficulty
  \item Idea density and conceptual difficulty
  \item Topic relevance to the reader
\end{parts}

\medskip




\noindent   \textbf{Question 58.} What is a key benefit of using the ``Think-Aloud'' strategy during reading?\pts{1}

\begin{parts}
  \item It verbalizes the reader's thought process (questions and predictions)
  \item It forces students to read faster
  \item It keeps the room quiet
  \item It replaces written tests
\end{parts}

\bigskip



\noindent   \textit{Reading Excerpt for Questions 59 to 61:}

\begin{quote}
   \textit{``The city was bustling with activity. People rushed to and fro, going about their daily business. The sounds of car horns, chatter, and construction filled the air.''}
\end{quote}
\smallskip




\noindent   \textbf{Question 59.} What is the atmosphere of the city in the passage?\pts{1}

\begin{parts}
  \item Bustling
  \item Peaceful
  \item Quiet
  \item Empty
\end{parts}

\medskip




\noindent   \textbf{Question 60.} What type of sounds are mentioned in the passage?\pts{1}

\begin{parts}
  \item A mix of urban sounds (horns, chatter, construction)
  \item Only natural sounds
  \item Silence
  \item Only animal sounds
\end{parts}

\medskip




\noindent   \textbf{Question 61.} What are people doing in the city?\pts{1}

\begin{parts}
  \item Going about their daily business
  \item Relaxing and resting
  \item Sleeping soundly
  \item Protesting
\end{parts}

\bigskip



\noindent   \textit{Reading Excerpt for Questions 62 to 64:}

\begin{quote}
   \textit{``The sun was setting over the ocean, painting the sky with hues of pink and orange. The waves gently lapped at the shore, creating a soothing melody.''}
\end{quote}
\smallskip




\noindent   \textbf{Question 62.} What is the tone of the passage?\pts{1}

\begin{parts}
  \item Calm and serene
  \item Sad and depressing
  \item Angry
  \item Humorous
\end{parts}

\medskip




\noindent   \textbf{Question 63.} What is happening to the sky in the passage?\pts{1}

\begin{parts}
  \item It is becoming colorful (pink and orange)
  \item It is getting darker
  \item It is getting lighter
  \item It is disappearing
\end{parts}

\medskip




\noindent   \textbf{Question 64.} What sound is created by the waves?\pts{1}

\begin{parts}
  \item A soothing melody
  \item A loud noise
  \item Complete silence
  \item A chaotic sound
\end{parts}

\bigskip



\noindent   \textit{Reading Excerpt for Questions 65 to 67:}

\begin{quote}
   \textit{``The library was a quiet haven for students. Rows of books lined the shelves, and the musty smell of old books filled the air.''}
\end{quote}
\smallskip




\noindent   \textbf{Question 65.} What is the atmosphere of the library?\pts{1}

\begin{parts}
  \item Quiet
  \item Noisy
  \item Crowded
  \item Empty
\end{parts}

\medskip




\noindent   \textbf{Question 66.} What fills the air in the library?\pts{1}

\begin{parts}
  \item The smell of old books
  \item The smell of food
  \item The sound of music
  \item The sound of chatter
\end{parts}

\medskip




\noindent   \textbf{Question 67.} For whom is the library a haven?\pts{1}

\begin{parts}
  \item Students
  \item Teachers
  \item Tourists
  \item Researchers
\end{parts}

\bigskip



\noindent   \textit{Reading Excerpt for Questions 68 to 70:}

\begin{quote}
   \textit{``The old woman sat on the bench, watching the world go by. She smiled as children played in the park, their laughter echoing through the air.''}
\end{quote}
\smallskip




\noindent   \textbf{Question 68.} What is the old woman doing?\pts{1}

\begin{parts}
  \item Watching the world go by
  \item Playing with children
  \item Watching TV
  \item Sleeping
\end{parts}

\medskip




\noindent   \textbf{Question 69.} What was her reaction to the children?\pts{1}

\begin{parts}
  \item Smiling
  \item Frowning
  \item Scowling
  \item Shuddering
\end{parts}

\medskip




\noindent   \textbf{Question 70.} What is echoing through the air?\pts{1}

\begin{parts}
  \item Children's laughter
  \item The sound of traffic
  \item The sound of construction
  \item Music
\end{parts}

\vfill
\rule{\linewidth}{0.4pt}

\begin{center}\small
\href{https://bhu-exam.vercel.app/}{Click here to visit our website}\\[4pt]
\href{https://me-aryan.vercel.app/}{Click here to visit our contributor}
\end{center}

\end{document}
